\documentclass[11pt]{article}
\usepackage[margin=1.2in]{geometry}
\usepackage{amsmath,amssymb,amsthm}
\usepackage{booktabs}
\usepackage{array,tabularx}
\usepackage[hidelinks]{hyperref}

\newcolumntype{Y}{>{\raggedright\arraybackslash}X}

\newtheorem{theorem}{Theorem}
\newtheorem{proposition}[theorem]{Proposition}
\newtheorem{lemma}[theorem]{Lemma}
\newtheorem{corollary}[theorem]{Corollary}
\theoremstyle{definition}
\newtheorem{remark}[theorem]{Remark}

\newcommand{\SB}{\mathrm{SB}}
\newcommand{\pr}{\mathbb{P}}
\newcommand{\E}{\mathbb{E}}

\title{Exact finite-frame inference under one-start systematic
probability-proportional-to-size sampling}
\author{Richie Sater\\
\small Independent researcher, Silver Spring, Maryland, USA\\
\small \texttt{richiesater@gmail.com}}
\date{September 2026}

\begin{document}
\maketitle

\begin{abstract}
The ordinary Stringer upper bound is routinely paired with monetary-unit
sampling, although a one-start systematic probability-proportional-to-size
(PPS) sample is dependent rather than i.i.d.  We determine what can be
guaranteed under the systematic design itself.  For every sample size
$n\ge2$ and every tail level $\alpha\in(0,1)$, a periodic two-item-per-cell
population makes binomial-factor Stringer undercover; the analogous
Poisson rule also fails whenever $-\log(\alpha)/n<1-\alpha$.  The systematic
sample mean nevertheless remains exactly design-unbiased.  This identity
implies both a sharp impossibility result and a remedy: every universally
valid taint-only bound must be at least $1-\alpha+\alpha y$ on the constant
sample $(y,\ldots,y)$, and that floor is attained.  Intersecting it with a
deterministic item-completion bound and then taking the maximum with capped
Stringer yields a design-valid overlay.  When the ordered frame is known,
exact inversion of the grid-mean randomization rank gives a pointwise
smaller least-favourable-completion bound, expressible as finitely many
rational linear programs.  For phase-partitioned frames this optimization
reduces to a minimum-coalition formula; equal phases give a closed form,
rational unequal phases admit an $O(KD)$ pseudo-polynomial algorithm, and
exact evaluation is weakly NP-hard.  Independent starts yield an explicit
all-zero formula.  In a 2,000-item, 100-draw example, exact inversion lowers
the valid fallback from $0.95$ to $0.90$.
\end{abstract}

\noindent\textbf{Keywords:} monetary-unit sampling; systematic PPS;
Stringer bound; finite-population inference; randomization inference;
linear programming; subset sum.

\section{Introduction}

Monetary-unit sampling selects book units with probability proportional to
recorded amount.  A common implementation places the ordered population on
the cumulative-book-value line, chooses one random start, and takes equally
spaced points thereafter.  This one-start systematic PPS design is easy to
implement, but the observations within its grid are dependent.  Professional
guidance permits systematic selection with random starts \cite{pcaob2315},
and statistical treatments of monetary-unit sampling emphasize the
importance of the actual unequal-probability design \cite{berger2021}.

For comparison, let $t_{(1)}\ge\cdots\ge t_{(n)}$ be observed taints in
nonincreasing order.  At tail level $\alpha$, the ordinary Stringer
calculation is
\begin{equation}\label{eq:sb}
 \SB_\alpha=p_n(0)+\sum_{j=1}^{n}
       \{p_n(j)-p_n(j-1)\}t_{(j)}.
\end{equation}
For the binomial convention, $p_n(j)$ is the upper Clopper--Pearson limit:
for $j<n$ it is the solution of
\[
 \sum_{i=0}^{j}\binom ni p^i(1-p)^{n-i}=\alpha,
 \qquad p_n(n)=1.
\]
For the Poisson convention, if $\lambda_j$ solves
\[
 e^{-\lambda_j}\sum_{i=0}^{j}\frac{\lambda_j^i}{i!}=\alpha,
\]
then $p_n(j)=\min(1,\lambda_j/n)$.  These factor sequences arise from
independent count models; applying them to one systematic grid therefore
requires a separate design argument.  Audit tables have long used the
Poisson convention \cite{aicpa2012}, while Stringer's original construction
uses binomial limits \cite{stringer1963}.

Prior work already warns against identifying systematic selection with
simple random sampling.  Hoogduin, Hall, Tsay, and Pierce \cite{hoogduin2015}
showed by simulation that this mismatch can materially distort assessed
risk.  Fienberg, Neter, and Leitch \cite{fienberg1977} obtained finite-sample
audit bounds by maximizing over a multinomial confidence region.  Their
discretized multinomial model differs from the fixed ordered-frame
randomization problem studied here, but it is an important
optimization-based predecessor. Recent work on randomized systematic
unequal-probability sampling derives exact second-order inclusion
probabilities for variance estimation \cite{vanberkel2026}; its target is
different from the one-sided taint-completion problem here. Thus the
contribution below is not a first warning about systematic sampling. It is an exact analysis: an
all-$n$, all-confidence binomial failure family, the sharp information limit
for taint-only bounds, a closed-form valid overlay, and a finite-frame
randomization inversion.

The main mechanism is the distinction between first moments and joint
dependence.  Averaging over the random start makes the systematic grid mean
exactly unbiased for the finite-population taint ratio.  A periodic ordering
can nevertheless make every point in the grid reveal the same phase.  The
first fact yields a Markov bound; the second proves that its value on
constant samples cannot be improved without using item identities or frame
structure.  Once the complete ordered frame is available, the unobserved
taints become nuisance variables in an exact randomization test.  Maximizing
the population ratio over compatible completions produces the sharper
finite-frame bound.

Sections~\ref{sec:one-start} and~\ref{sec:finite-frame} develop these two
levels.  Section~\ref{sec:disjoint} identifies the phase-partitioned case in
which the general linear programs collapse to subset sum and derives the
independent-start formula.  Section~\ref{sec:calculation-boundary} states the
role of the exact finite computations.  The universal theorems assume a
fixed ordered population with positive known book weights, taints in
$[0,1]$, and a uniform random start; they do not cover successive PPS,
stratification, nonrandom starts, uncertain book amounts, or negative
taints.

\section{One-start design: unbiasedness, failure, and a sharp safeguard}
\label{sec:one-start}

The key split is between first-moment calibration and dependence: the sample
average is design-unbiased, but the complete grid can behave as a single
random observation. We first prove that separation, then derive the resulting
failure family and use the same first-moment identity to construct a
closed-form valid reporting rule.

\subsection{Finite-population model and the two-phase construction}

Let an ordered finite population have positive recorded amounts
$w_1,\ldots,w_N$ and taints $t_1,\ldots,t_N\in[0,1]$. Put
\[
 W=\sum_{i=1}^Nw_i,
 \qquad
 \theta=\frac1W\sum_{i=1}^Nw_it_i.
\]
Place the items consecutively on $[0,W)$ and let $\tau(u)$ equal the taint
of the item containing $u$. For intended sample size $n$, write $h=W/n$,
draw $R\sim\operatorname{Unif}[0,h)$, and observe
\[
 T_k=\tau(R+kh),\qquad k=0,\ldots,n-1.
\]
Endpoints have probability zero and may be assigned by any fixed convention.
This is the continuous book-unit form of one-start systematic PPS; the
integer-book-unit version chooses one of the positions in an integer interval.

\begin{proposition}[Design-unbiased systematic mean]
\label{prop:sysunbiased}
For every such ordered population,
\[
 \E_R\overline T=\theta,
 \qquad
 \overline T=\frac1n\sum_{k=0}^{n-1}T_k.
\]
\end{proposition}

\begin{proof}
Integration over the systematic cells gives
\[
 \E_R\overline T
 =\frac1n\sum_{k=0}^{n-1}\frac1h
   \int_0^h\tau(r+kh)\,dr
 =\frac1W\int_0^W\tau(u)\,du
 =\theta.
\]
\end{proof}

Unbiasedness does not imply a binomial count law. Fix $0\le y<1$ and
$0<q<1$. In each of the $n$ systematic cells put first an item of weight
$qh$ and taint one, then an item of weight $(1-q)h$ and taint $y$. The
complete sample is
\begin{equation}\label{eq:systematictwophase}
 (T_0,\ldots,T_{n-1})=
 \begin{cases}
 (1,\ldots,1),&\text{with probability }q,\\
 (y,\ldots,y),&\text{with probability }1-q,
 \end{cases}
 \qquad
 \theta=y+q(1-y).
\end{equation}
All items are smaller than one systematic interval, so this construction
does not depend on duplicate-hit or certainty-item conventions.

\begin{theorem}[Exact systematic-PPS failure family]
\label{thm:systematicfailure}
For every $n\ge2$ and every $\alpha\in(0,1)$, some ordered finite population
of the form \eqref{eq:systematictwophase} gives ordinary binomial Stringer
coverage strictly below $1-\alpha$. For ordinary Poisson factors the same
conclusion holds whenever
\begin{equation}\label{eq:poissonsystematiccondition}
 \frac{-\log\alpha}{n}<1-\alpha.
\end{equation}
Thus Poisson Stringer fails on some population at 90\%, 95\%, and 99\%
confidence for every $n\ge3,4,5$, respectively. Final capping at one does
not alter these conclusions.
\end{theorem}

\begin{proof}
For binomial factors, $p_n(0)=1-\alpha^{1/n}$ and telescoping the increments
on the low phase gives
\[
 \SB_{\rm B}(y,\ldots,y)
 =p_n(0)+\{1-p_n(0)\}y
 =1-\alpha^{1/n}(1-y).
\]
Because $\alpha^{1/n}>\alpha$, choose
\[
 1-\alpha^{1/n}<q<1-\alpha.
\]
Then the low-phase bound is below $\theta$, while that phase has probability
$1-q>\alpha$. This proves binomial undercoverage. The high-phase binomial
bound is one, so coverage in this family is exactly $q$.

For Poisson factors take $y=0$. The all-zero bound is
$\lambda_0/n=(-\log\alpha)/n$. Under
\eqref{eq:poissonsystematiccondition}, choose
$(-\log\alpha)/n<q<1-\alpha$. Again the low phase is a noncoverage event of
probability greater than $\alpha$. The stated integer thresholds follow by
direct substitution. In both arguments the open interval for $q$ contains a
rational $a/m$; taking systematic interval $h=m$ gives the same construction
with integer book amounts $a$ and $m-a$ in every cell.
\end{proof}

The smallest exact witness in a unit-item census at 95\% confidence has only
nine book units. With $n=3$, interval three, and ordered taints
\[
 (1,1,0\mid1,1,0\mid1,1,0),
\]
the three starts return three ones, three ones, and three zeros. The target is
$2/3$, whereas $p_3(0)=1-0.05^{1/3}<2/3$, so exact coverage is $2/3$ rather
than $0.95$. Exact enumeration of all $2{,}188$ ordered binary populations
and $4{,}444$ phases through total size nine finds no smaller unit-item
witness; this finite minimality statement is limited to that enumerated class.

A practical-size version makes the failure scale explicit. Take total
recorded amount 1,000, $n=100$, and interval 10. Repeat 100 times an item of
amount one and taint one followed by an item of amount nine and taint zero.
One start gives 100 ones and the other nine starts give 100 zeros, while
$\theta=0.1$. On the high phase the binomial bound is one. For Poisson
factors, the exact inequality
$\pr\{\operatorname{Pois}(10)\le10\}>0.10$ implies
$\lambda_{100}>10$, so that phase also covers at all three levels below.
Consequently, both ordinary factor conventions have exact coverage 0.1.

\begin{table}[htbp]
\centering
\small
\begin{tabular}{@{}rrrrrrr@{}}
\toprule
confidence & B zero & P zero & ordinary coverage & safe output & B uplift & P uplift\\
\midrule
$90\%$ & 0.0227628 & 0.0230259 & 0.1000000 & 0.1000000 & 0.0772372 & 0.0769741\\
$95\%$ & 0.0295130 & 0.0299573 & 0.1000000 & 0.1000000 & 0.0704870 & 0.0700427\\
$99\%$ & 0.0450074 & 0.0460517 & 0.1000000 & 0.1000000 & 0.0549926 & 0.0539483\\
\bottomrule
\end{tabular}
\caption{The 100-draw periodic population. ``B zero'' and ``P zero'' are the
binomial- and Poisson-factor all-zero outputs. The design-valid safe output is
defined below. Displays are rounded from exact rational brackets in the
systematic-PPS calculation.}
\label{tab:systematicexample}
\end{table}

\subsection{A sharp taint-only floor and an item-aware improvement}

The construction also identifies the unavoidable price of robustness if a
rule sees only the taint vector. Item identities and book weights will then
provide a separate, strictly useful information channel.

\begin{proposition}[Necessary taint-only diagonal]
\label{prop:systematicfloor}
If $V_\alpha:[0,1]^n\to[0,1]$ has systematic-design coverage at least
$1-\alpha$ for every ordered finite population and depends only on the
observed taints, then, for every $0\le y<1$,
\[
 V_\alpha(y,\ldots,y)\ge1-\alpha+\alpha y.
\]
\end{proposition}

\begin{proof}
If the inequality failed, choose
\[
 \frac{V_\alpha(y,\ldots,y)-y}{1-y}<q<1-\alpha
\]
and use \eqref{eq:systematictwophase}. Its target
$y+q(1-y)$ would exceed the reported low-phase value, and that phase would
have probability $1-q>\alpha$.
\end{proof}

The lower bound is attained on every diagonal by
\begin{equation}\label{eq:systematicmarkov}
 M_\alpha=1-\alpha+\alpha\overline T.
\end{equation}
Indeed, with $X=1-\overline T$, Proposition~\ref{prop:sysunbiased} gives
$\E_R X=1-\theta$. The event $M_\alpha<\theta$ is
$X>(\E_R X)/\alpha$, whose probability is at most $\alpha$ by Markov's
inequality. Thus one random start can carry only one observation's worth of
distribution-free information, regardless of the number of points in its
grid.

For a sampled set $S$ of distinct item identities, define the deterministic
completion bound
\[
 C(S)=\frac1W\left(\sum_{i\in S}w_it_i+\sum_{i\notin S}w_i\right).
\]
Because each unobserved taint is at most one, $C(S)\ge\theta$ pointwise.
This extra information permits a smaller valid floor.

\begin{proposition}[Systematic-safe Stringer overlay]
\label{prop:systematicsafeguard}
Put
\begin{equation*}
 H_\alpha=\min\{M_\alpha,C(S)\},
 \qquad
 \SB_\alpha^{\rm sys}
 =\max\{\min(1,\SB_\alpha),H_\alpha\}.
\tag{SYS}\label{eq:systematicsafe}
\end{equation*}
For either the binomial- or Poisson-factor ordinary calculation,
$\SB_\alpha^{\rm sys}$ has coverage at least $1-\alpha$ under the one-start
systematic design for every ordered finite population. It is never below
the ordinary output after its natural cap at one.
\end{proposition}

\begin{proof}
Whenever $M_\alpha\ge\theta$, the deterministic inequality
$C(S)\ge\theta$ implies $H_\alpha\ge\theta$. Hence $H_\alpha$ has
noncoverage at most $\alpha$. The maximum in \eqref{eq:systematicsafe} is at
least $H_\alpha$, which proves coverage, and at least the capped ordinary
output by definition.
\end{proof}

In the all-zero phase of Table~\ref{tab:systematicexample}, the 100 audited
taint-zero items comprise 900 of the 1,000 recorded units. At 95\% confidence,
$M_{0.05}=0.95$ but $C(S)=0.10$, so $H_{0.05}=0.10$. The safe output is
therefore exactly 0.10, an uplift of about 0.070487 over binomial Stringer and
0.070043 over Poisson Stringer, rather than the 0.92 uplift of the taint-only
floor. On the all-one phase, $H_\alpha=1$.

\section{Exact finite-frame minimax inversion}
\label{sec:finite-frame}

The closed-form bound uses the observed identities only through $C(S)$. A
fixed known frame contains more information: it specifies which items every
other random start would have selected. We now invert that complete finite
randomization distribution.

Combine random-start intervals having the same item-hit multiplicities into
phase atoms $a=1,\ldots,K$, with probabilities $\pi_a>0$. Let $m_{ai}$ be
the number of hits on item $i$ in atom $a$, and for a complete candidate
taint vector $u\in[0,1]^N$ put
\[
 Y_a(u)=\frac1n\sum_{i=1}^N m_{ai}u_i,
 \qquad
 \theta(u)=\frac1W\sum_{i=1}^Nw_iu_i.
\]
Suppose the observed phase audits the distinct set $S$, returns values
$x=(x_i)_{i\in S}$, and has grid mean $y$. Define the exact lower-tail rank
\[
 F_u(y)=\sum_{a=1}^K\pi_a\mathbf 1\{Y_a(u)\le y\}
\]
and the least-favourable-completion bound
\begin{equation*}\label{eq:finiteframebound}
 B^{\rm ff}_\alpha(S,x)
 =\max\left(
 \left\{\theta(u):0\le u\le1,\ u_S=x,\ F_u(y)>\alpha\right\}
 \cup\{0\}
 \right).
\tag{FF}
\end{equation*}
Here ``minimax'' refers only to maximizing the target over nuisance-taint
completions for this pre-specified grid-mean ordering. It does not assert
optimality among all possible test orderings or confidence procedures.

\begin{proposition}[Exact finite-frame inversion]
\label{prop:systematicminimax}
For every fixed ordered frame, $B^{\rm ff}_\alpha$ has one-start systematic
design coverage at least $1-\alpha$. Pointwise,
\begin{equation}\label{eq:finiteframedominance}
 B^{\rm ff}_\alpha(S,x)
 \le \min\{M_\alpha,C(S)\}=H_\alpha.
\end{equation}
Moreover, for rational book weights and observed taints, the value in
\eqref{eq:finiteframebound} is the maximum of finitely many rational linear
programs indexed by inclusion-minimal phase coalitions of probability
strictly greater than $\alpha$.
\end{proposition}

\begin{proof}
Fix the true population $t$, and let $F_t$ be the CDF of its grid mean over
the random start. If $B^{\rm ff}_\alpha(S_R,t_{S_R})<\theta(t)$, then the true
completion is absent from the feasible set in \eqref{eq:finiteframebound};
therefore $F_t(Y_R(t))\le\alpha$. For any discrete random variable $Y$,
\[
 \Pr\{F_Y(Y)\le\alpha\}\le\alpha,
\]
because the event, when nonempty, has probability equal to the CDF at its
largest included support point. This proves design coverage.

Every compatible completion is bounded by $C(S)$. For the other inequality,
take a feasible $u$ and view $Y$ as the random phase mean $Y_a(u)$. By
Proposition~\ref{prop:sysunbiased}, $\E Y=\theta(u)$. If $y<1$, Markov's
inequality applied to $1-Y$ gives
\[
 \alpha<F_u(y)=\Pr\{1-Y\ge1-y\}
 \le\frac{1-\theta(u)}{1-y},
\]
so $\theta(u)<1-\alpha+\alpha y=M_\alpha$. If $y=1$, the inequality follows
from the unit range. The feasible set is a finite union of compact polytopes,
so its maximum is attained; this proves \eqref{eq:finiteframedominance}.

Finally, for a phase coalition $A$ with
$\sum_{a\in A}\pi_a>\alpha$, impose the linear constraints
\[
 0\le u_i\le1,\qquad u_i=x_i\ (i\in S),\qquad
 \sum_i m_{ai}u_i\le ny\ (a\in A),
\]
and maximize the linear objective $\theta(u)$. The condition $F_u(y)>
\alpha$ holds exactly when the low-mean phase set contains such a coalition.
Every coalition of mass above $\alpha$ contains an inclusion-minimal one, and
deleting constraints cannot lower a maximum. Taking the maximum of these
rational LP values therefore gives exactly \eqref{eq:finiteframebound}.
\end{proof}

Thus the sharper pre-specified overlay
\begin{equation*}\label{eq:finiteframesafe}
 \SB_{\alpha}^{\rm ff-sys}
 =\max\{\min(1,\SB_\alpha),B^{\rm ff}_\alpha\}
\tag{FF-SYS}
\end{equation*}
is design-valid, never below capped ordinary Stringer, and never above the
earlier overlay \eqref{eq:systematicsafe}. For rational inputs, the finite maximum can be evaluated exactly by rational
linear programming; equality of feasible primal and dual objectives certifies
each optimum.

\section{Phase partitions, exact algorithms, and independent starts}
\label{sec:disjoint}

The general LP representation simplifies when random starts see separate
parts of the frame. Call a frame \emph{phase-partitioned} when the supports of
the distinct phase-incidence atoms are pairwise disjoint and cover all items.
Design unbiasedness then has the decomposition
\[
 \theta(u)=\sum_{a=1}^K\pi_aY_a(u),
\]
and the unobserved phase means can be varied independently.

\begin{proposition}[Unequal- and equal-phase formulas]
\label{prop:disjointphase}
Suppose phase $j$ is observed with mean $y$ on a phase-partitioned frame. Put
\[
 \rho_j(\alpha)
 =\min\{\pi(A):j\in A,\ \pi(A)>\alpha\},
 \qquad \pi(A)=\sum_{a\in A}\pi_a.
\]
Then
\begin{equation}\label{eq:disjointphase}
 B^{\rm ff}_\alpha(j,y)=1-(1-y)\rho_j(\alpha).
\end{equation}
If there are $m$ equiprobable phases, this becomes
\begin{equation}\label{eq:equalphase}
 B^{\rm ff}_\alpha(j,y)
 =1-\frac{\lfloor m\alpha\rfloor+1}{m}(1-y).
\end{equation}
\end{proposition}

\begin{proof}
For any feasible completion let $A=\{a:Y_a(u)\le y\}$. It contains the
observed phase and has $\pi(A)>\alpha$. Therefore
\[
 \theta(u)
 \le y\pi(A)+1-\pi(A)
 =1-(1-y)\pi(A)
 \le1-(1-y)\rho_j(\alpha).
\]
Choose a minimizing coalition $A_\star$. Keep the observed taints fixed, set
every taint in each other phase of $A_\star$ equal to $y$, and set every
taint outside $A_\star$ equal to one. The resulting phase means are $y$ on
$A_\star$ and one elsewhere, so the completion is feasible and attains the
displayed upper bound. This proves \eqref{eq:disjointphase}. With equal
phase probabilities, the least cardinality having mass strictly above
$\alpha$ is $\lfloor m\alpha\rfloor+1$, proving \eqref{eq:equalphase}.
\end{proof}

For rational unequal probabilities, the minimum in
\eqref{eq:disjointphase} is a minimum-excess subset sum. The next result gives
both an exact algorithm and the obstruction to a bit-polynomial algorithm.

\begin{proposition}[Exact dynamic program and weak NP-hardness]
\label{prop:disjointdp}
Write $\pi_a=c_a/D$ with positive integers $c_a$ summing to $D$. There is an
exact algorithm for \eqref{eq:disjointphase} using $O(KD)$ time and $O(D)$
memory. Under binary encoding, exact evaluation is weakly NP-hard even for
$n=1$ and an observed taint of zero.
\end{proposition}

\begin{proof}
The smallest integer coalition mass strictly above $\alpha$ is
\[
 h=\lfloor\alpha D\rfloor+1.
\]
Starting with the compulsory mass $c_j$, list the other phase masses as
$d_1,\ldots,d_{K-1}$ and form the reachable sets
\[
 R_0=\{0\},\qquad R_{\ell+1}=R_\ell\cup(R_\ell+d_{\ell+1}).
\]
Every optional coalition yields one reachable sum and conversely. Hence the
first reachable $s$ satisfying $c_j+s\ge h$ gives
$\rho_j(\alpha)=(c_j+s)/D$. A Boolean array with back-pointers implements
the recurrence in the stated time and memory.

For hardness, reduce from \textsc{Subset Sum} \cite{gareyjohnson1979}. Given
positive integers $a_1,\ldots,a_k$ and target
$1\le T\le A:=\sum_i a_i$, take a one-draw systematic-PPS frame with weights
\[
 w_0=1,\qquad w_i=2a_i,\qquad W=1+2A,
\]
set $\alpha=2T/W$, and observe item $0$ with taint zero. Each item is then a
disjoint phase of probability $w_i/W$. A zero-taint coalition containing
item $0$ has weight $1+2\sum_{i\in I}a_i$ and probability strictly above
$\alpha$ exactly when $\sum_{i\in I}a_i\ge T$. If
\[
 s_\star=\min\{\textstyle\sum_{i\in I}a_i:
                     \sum_{i\in I}a_i\ge T\},
\]
then Proposition~\ref{prop:disjointphase} gives
\[
 B^{\rm ff}_\alpha=1-\frac{1+2s_\star}{W}.
\]
Thus the original instance has a subset summing exactly to $T$ if and only if
the exact bound equals $1-(1+2T)/W$. The reduction has polynomial encoding
length. Together with the pseudo-polynomial dynamic program, this proves the
claim.
\end{proof}

Formula \eqref{eq:equalphase} gives a practical-size strict improvement
without solving an LP. Take 2,000 unit items,
$n=100$, and systematic interval 20. The 20 equiprobable phases select
disjoint sets of 100 items. After one all-zero phase at 95\% confidence,
\[
 M_{0.05}=C(S)=H_{0.05}=0.95.
\]
Here $\lfloor20(0.05)\rfloor+1=2$, so
\begin{equation}\label{eq:finiteframebenchmark}
 B^{\rm ff}_{0.05}=\frac{1800}{2000}=0.90.
\end{equation}
The general LP solver independently enumerates all
$\binom{20}{2}=190$ minimal coalitions and returns the same witness with
lower-tail mass $0.10$; the specialized dynamic program returns the same
minimum mass directly. Both ordinary all-zero
Stringer outputs are below 0.03, so the finite-frame overlay reports 0.90
under either factor convention, five percentage points below the previous
valid fallback.

\subsection{Independent starts}

The disjoint equal-phase case also gives an exact replication law rather than
only a generic bounded-mean reduction.

\begin{proposition}[Independent-start all-zero formula]
\label{prop:multistartzero}
Use $r\ge1$ independent starts on $m$ equiprobable disjoint phases, and
invert the lower-tail rank of the mean of the $r$ grid averages. Suppose all
observed taints are zero and the starts visit $d$ distinct phases. Put
\[
 q_{m,r,\alpha,d}
 =\max\{d,\lfloor m\alpha^{1/r}\rfloor+1\}.
\]
Then
\begin{equation}\label{eq:multistartzero}
 B^{\rm ff}_{\alpha,r}=1-\frac{q_{m,r,\alpha,d}}m.
\end{equation}
\end{proposition}

\begin{proof}
For a candidate completion, let $Z$ be the set of zero-mean phases. Since
taints are nonnegative, the mean of $r$ independently selected phase means is
at most its observed value zero exactly when every selected phase lies in
$Z$. Its lower-tail probability is therefore $(|Z|/m)^r$. Compatibility
forces all $d$ observed phases into $Z$, and feasibility requires this
probability to exceed $\alpha$. Thus $|Z|\ge q_{m,r,\alpha,d}$. Such phases
contribute zero to the target and the others contribute at most $1/m$ each,
which proves the upper bound. Assigning taint zero to exactly
$q_{m,r,\alpha,d}$ phases containing the observed ones, and taint one to all
remaining phases, attains equality.
\end{proof}

For $m=20$ and $\alpha=0.05$, distinct all-zero starts give exact bounds
$0.90$, $0.75$, and $0.60$ for $r=1,2,3$. For unequal rational phases, the
same dynamic program replaces its one-start threshold by the least integer
$h$ satisfying $(h/D)^r>\alpha$ and makes every distinct observed phase
compulsory.

If $r$ independent random starts are used, their $r$ grid averages are i.i.d.
variables in $[0,1]$ with mean $\theta$. Any valid distribution-free
bounded-mean procedure may be applied to those averages. The effective
replication is the number of independent starts, not the number of points
within one grid.

\section{Exact calculations and mathematical scope}
\label{sec:calculation-boundary}

The universal failure family, unbiasedness identity, sharp taint-only floor,
valid overlays, finite-frame inversion, phase-partition formulas, complexity
result, and independent-start formula are ordinary mathematical proofs.  No
finite experiment is extrapolated to establish those statements.

Exact rational programs establish the finite claims reported above.  The
nine-unit calculation enumerates all $2{,}188$ ordered binary unit-item
populations and all $4{,}444$ associated phases in its stated search class.
The 100-draw calculation uses rational phase probabilities and rigorous
brackets for the binomial and Poisson factors.  For the 2,000-item example,
the general finite-frame solver enumerates the 190 inclusion-minimal
coalitions and certifies each rational LP optimum by matching feasible primal
and dual objectives; a separate denominator-cleared subset-sum program gives
the same minimum coalition mass.  These computations certify only the
displayed census and examples.  The programs use exact rational arithmetic
apart from outward enclosures of transcendental confidence factors.

The sampling-design scope is also exact.  One-start results average over one
uniform start on a fixed ordered frame.  The multiple-start proposition
changes the design to independent starts and, in its closed form, assumes
equiprobable disjoint phases and an all-zero observation.  Alternative test
orderings, overlapping-phase multiple-start formulas, stratification,
successive PPS, nonrandom starts, uncertain book amounts, negative taints,
and different audit estimands require separate arguments.

\section{Discussion}

The first lesson is that unbiasedness and replication are different
properties.  Averaging over the random start recovers the finite-population
taint ratio, yet a periodic ordering can make all $n$ observed taints encode
one phase.  Consequently no i.i.d. coverage theorem for the ordinary
Stringer calculation could by itself validate its use under one-start
systematic PPS.  The reusable statistical object is the mean of each
independently randomized grid.

The second lesson is an exact value-of-information statement.  If the rule
sees only the taint vector, every universally design-valid procedure must pay
the one-observation diagonal floor $1-\alpha+\alpha y$.  Sampled item
identities and book weights can lower that floor because they determine how
much of the total recorded amount has already been resolved.  The complete
ordered frame contributes still more: it turns every possible start into a
constraint on nuisance-taint completions.  Randomization-rank inversion
combines these constraints and is pointwise no larger than both closed-form
fallbacks.

The phase-partition formula isolates the combinatorial core of this
inversion.  Disjoint phase supports make the target a weighted average of
independently adjustable phase means, so inference becomes the problem of
finding the least phase mass strictly above $\alpha$.  Equal phases expose
the discrete correction in closed form; unequal rational phases connect the
exact bound to subset sum, explaining both the pseudo-polynomial algorithm
and the weak NP-hardness boundary.  With independent starts, the required
zero-phase mass rises from $\alpha$ to $\alpha^{1/r}$.  These formulas make
the inferential gain from additional random starts quantitative rather than
merely qualitative.

\section*{Data and code availability}

Source code and the exact calculation data used in the finite examples are
available in the project repository at
\href{https://github.com/RichieSater/stringer-bound}
{github.com/RichieSater/stringer-bound}.

\section*{Funding}

The author received no external funding for this work.

\section*{Competing interests}

The author declares no competing interests.

\begin{thebibliography}{99}

\bibitem{aicpa2012} American Institute of Certified Public Accountants,
\emph{Audit Sampling}, AICPA Audit Guide, 2012.

\bibitem{berger2021} Y.~G. Berger, P.~M. Chiodini and M.~Zenga, Bounds for
monetary-unit sampling in auditing: an adjusted empirical likelihood
approach, \emph{Statistical Papers} \textbf{62} (2021), 2739--2761,
\url{https://doi.org/10.1007/s00362-020-01209-w}.

\bibitem{fienberg1977} S.~E. Fienberg, J.~Neter and R.~A. Leitch,
Estimating the total overstatement error in accounting populations,
\emph{Journal of the American Statistical Association} \textbf{72} (1977),
295--302, \url{https://doi.org/10.1080/01621459.1977.10480993}.

\bibitem{gareyjohnson1979} M.~R. Garey and D.~S. Johnson,
\emph{Computers and Intractability: A Guide to the Theory of
NP-Completeness}, W.~H. Freeman, San Francisco, 1979.

\bibitem{hoogduin2015} E.~Hoogduin, T.~W. Hall, B.-Y. Tsay and B.~Pierce,
Does systematic selection lead to unreliable risk assessments in
monetary-unit sampling applications?, \emph{Auditing: A Journal of Practice
\& Theory} \textbf{34} (2015), no.~4, 85--107,
\url{https://doi.org/10.2308/ajpt-51081}.

\bibitem{pcaob2315} Public Company Accounting Oversight Board, AS~2315:
Audit sampling,
\href{https://pcaobus.org/oversight/standards/auditing-standards/details/AS2315}
{PCAOB standard page} (accessed August 31, 2026).

\bibitem{stringer1963} K.~W. Stringer, Practical aspects of statistical
sampling in auditing, in \emph{Proceedings of the Business and Economic
Statistics Section}, American Statistical Association, 1963, pp.~405--411.

\bibitem{vanberkel2026} K.~van Berkel and E.~Vondenhoff, Exact and
approximation formulas for second-order inclusion probabilities in randomized
systematic sampling with unequal probabilities and without replacement,
\emph{Survey Methodology} \textbf{52} (2026), no.~1, 235--258,
\url{https://www150.statcan.gc.ca/n1/pub/12-001-x/2026001/article/00003-eng.htm}.

\end{thebibliography}

\end{document}
